\documentclass[10pt,a4paper]{amsart}
\usepackage{amsmath,amssymb,amsthm,graphicx}
% Retained preamble for the Stacks Project chapter {CH} (shared project preamble).
% The chapter file pinned-original-source.tex only carries \input{{preamble}} on its first line,
% so no preamble text lives inside the chapter itself.  The Stacks Project uses one single shared
% preamble for every chapter: the snapshot stages it three times (sources/stacks-chow/preamble.tex,
% sources/stacks-etale-cohomology/preamble.tex and focused-seeds/AGP/sources/stacks-snapshot-20261011/
% preamble.tex) and those copies are identical in size.  The block below is copied mechanically,
% line for line, from sources/stacks-chow/preamble.tex lines 14-18 and 154-217: the Xy-pic setup and
% the native theorem declarations together with the mathematical macros.  Nothing was added, removed
% or redefined; the wrapper class and the amsmath/amssymb/amsthm/graphicx base packages are supplied
% by the assembler.
% For commutative diagrams we use Xy-pic
% Chapter: stacks-cohomology.tex -- Cohomology of Sheaves.
\usepackage[all]{xy}

% We use 2cell for 2-commutative diagrams.
\xyoption{2cell}

% Theorem environments.
%
\theoremstyle{plain}
\newtheorem{theorem}[subsection]{Theorem}
\newtheorem{proposition}[subsection]{Proposition}
\newtheorem{lemma}[subsection]{Lemma}

\theoremstyle{definition}
\newtheorem{definition}[subsection]{Definition}
\newtheorem{example}[subsection]{Example}
\newtheorem{exercise}[subsection]{Exercise}
\newtheorem{situation}[subsection]{Situation}

\theoremstyle{remark}
\newtheorem{remark}[subsection]{Remark}
\newtheorem{remarks}[subsection]{Remarks}

\numberwithin{equation}{subsection}

% Macros
%
\def\lim{\mathop{\mathrm{lim}}\nolimits}
\def\colim{\mathop{\mathrm{colim}}\nolimits}
\def\Spec{\mathop{\mathrm{Spec}}}
\def\Hom{\mathop{\mathrm{Hom}}\nolimits}
\def\Ext{\mathop{\mathrm{Ext}}\nolimits}
\def\SheafHom{\mathop{\mathcal{H}\!\mathit{om}}\nolimits}
\def\SheafExt{\mathop{\mathcal{E}\!\mathit{xt}}\nolimits}
\def\Sch{\mathit{Sch}}
\def\Mor{\mathop{\mathrm{Mor}}\nolimits}
\def\Ob{\mathop{\mathrm{Ob}}\nolimits}
\def\Sh{\mathop{\mathit{Sh}}\nolimits}
\def\NL{\mathop{N\!L}\nolimits}
\def\CH{\mathop{\mathrm{CH}}\nolimits}
\def\proetale{{pro\text{-}\acute{e}tale}}
\def\etale{{\acute{e}tale}}
\def\QCoh{\mathit{QCoh}}
\def\Ker{\mathop{\mathrm{Ker}}}
\def\Im{\mathop{\mathrm{Im}}}
\def\Coker{\mathop{\mathrm{Coker}}}
\def\Coim{\mathop{\mathrm{Coim}}}

% Boxtimes
%
\DeclareMathSymbol{\boxtimes}{\mathbin}{AMSa}{"02}

%
% Macros for moduli stacks/spaces
%
\def\QCohstack{\mathcal{QC}\!\mathit{oh}}
\def\Cohstack{\mathcal{C}\!\mathit{oh}}
\def\Spacesstack{\mathcal{S}\!\mathit{paces}}
\def\Quotfunctor{\mathrm{Quot}}
\def\Hilbfunctor{\mathrm{Hilb}}
\def\Curvesstack{\mathcal{C}\!\mathit{urves}}
\def\Polarizedstack{\mathcal{P}\!\mathit{olarized}}
\def\Complexesstack{\mathcal{C}\!\mathit{omplexes}}
% \Pic is the operator that assigns to X its picard group, usage \Pic(X)
% \Picardstack_{X/B} denotes the Picard stack of X over B
% \Picardfunctor_{X/B} denotes the Picard functor of X over B
\def\Pic{\mathop{\mathrm{Pic}}\nolimits}
\def\Picardstack{\mathcal{P}\!\mathit{ic}}
\def\Picardfunctor{\mathrm{Pic}}
\def\Deformationcategory{\mathcal{D}\!\mathit{ef}}
\title{Canonical map from a complex of modules to the derived pushforward, functorial in the triple}
\author{The Stacks Project Authors}
\date{stacks-project:cohomology.tex:local-snapshot-20261011; GNU FDL 1.2 or later}
\begin{document}
\maketitle
\noindent\textit{Editorial note.} All delivered mathematical text is the author's own, copied from the pinned Stacks Project chapter file cohomology.tex; no statement or proof step was rewritten, reordered or removed.\par\medskip
\noindent\textit{Editorial note.} The 'moreover' clause on functoriality is discharged by the author through the construction itself; no separate argument is supplied in the source and none was added.\par\medskip
\noindent\textit{Editorial note.} The preamble is the chapter's retained minimal preamble (sources/stacks-cohomology/retained-preamble.tex lines 1-81), a verbatim excerpt of the native preamble: Xy-pic setup plus the native theorem declarations and mathematical macros.\par\medskip
\noindent\textit{Editorial note.} The label lemma-functoriality is defined only in cohomology.tex, so no cross-chapter restatement had to be weighed for this card.\par\medskip
\noindent\textit{Allowed standard prerequisites (editorial).}\begin{itemize}
\item Injective resolutions in the category of O\_X-modules. Same chapter, the existence of enough injectives among O\_X-modules, used implicitly at proof line 2348. Supplies the resolution F to I from which the representative f\_*I of Rf\_*F is built.
\item The derived category D\textasciicircum{}+(Y) as a localization of K\textasciicircum{}+(Y). Homology, the construction of D(A) and the localization functor j\_Y: K\textasciicircum{}+(Y) to D\textasciicircum{}+(Y) named at proof line 2357. Supplies the last step of the construction and the meaning of a canonical morphism in D\textasciicircum{}+(Y).
\item Definition of the derived pushforward Rf\_*. Same chapter, Rf\_* computed on an injective resolution, used at proof lines 2349-2350. Supplies that f\_*I represents Rf\_*F in K\textasciicircum{}+(O\_Y).
\end{itemize}\medskip
\begin{lemma}
\label{lemma-functoriality}
Let $f : X \to Y$ be a morphism of ringed spaces.
Let $\mathcal{G}^\bullet$, resp.\ $\mathcal{F}^\bullet$ be
a bounded below complex of $\mathcal{O}_Y$-modules,
resp.\ $\mathcal{O}_X$-modules. Let
$\varphi : \mathcal{G}^\bullet \to f_*\mathcal{F}^\bullet$
be a morphism of complexes. There is a canonical morphism
$$
\mathcal{G}^\bullet
\longrightarrow
Rf_*(\mathcal{F}^\bullet)
$$
in $D^{+}(Y)$. Moreover this construction is functorial in the triple
$(\mathcal{G}^\bullet, \mathcal{F}^\bullet, \varphi)$.
\end{lemma}

\begin{proof}
Choose an injective resolution $\mathcal{F}^\bullet \to \mathcal{I}^\bullet$.
By definition $Rf_*(\mathcal{F}^\bullet)$ is represented by
$f_*\mathcal{I}^\bullet$ in $K^{+}(\mathcal{O}_Y)$.
The composition
$$
\mathcal{G}^\bullet \to f_*\mathcal{F}^\bullet \to f_*\mathcal{I}^\bullet
$$
is a morphism in $K^{+}(Y)$ which turns
into the morphism of the lemma upon applying the
localization functor $j_Y : K^{+}(Y) \to D^{+}(Y)$.
\end{proof}

\end{document}
